% ---- Method -----------------------------------------------------------------
% Target length: ~1 page.
% Structure: SVG representation, base model, training formulation, LoRA, curriculum.

\subsection{SVG as Code}

SVG is an XML-based language for describing two-dimensional vector graphics.
An icon SVG typically consists of a root \texttt{<svg>} element with a
\texttt{viewBox} attribute defining the coordinate space, followed by one or
more geometric primitives: \texttt{<path>} elements with \texttt{d} attributes
encoding sequences of drawing commands (MoveTo, LineTo, CubicBézier,
EllipticalArc, ClosePath), and optionally \texttt{<circle>}, \texttt{<rect>},
\texttt{<polygon>}, or \texttt{<g>} grouping elements.

This structure is naturally representable as a text sequence, and modern code
language models pre-trained on source code corpora have already seen XML and
SVG-like content in their training data. We leverage this prior knowledge by
treating icon generation as instruction-following code generation: given a
natural language description as the user instruction, produce a complete,
valid SVG document as the assistant response.

\subsection{Base Model}

We use Qwen3.5-9B \citep{qwen3} as the base model. Qwen3.5 represents the
current generation of Qwen language models \citep{qwen25}, featuring a 128K
token context window (important for complex SVGs), strong performance on code
generation benchmarks, and native instruction-following capabilities through
the ChatML prompt format. The 9B parameter scale balances generation quality
with the memory and compute constraints of local Apple Silicon hardware.

\subsection{Instruction-Tuning Formulation}

Each training example is formatted as a three-turn conversation in the
ChatML format used by Qwen3.5:

\begin{lstlisting}[style=svgstyle, caption={Training record format (ChatML)}]
{"messages": [
  {"role": "system",
   "content": "You are an SVG icon generator. Given a description,
               output clean, valid SVG code for a single icon. Use
               currentColor for fill and stroke. Output only the
               SVG element, no explanation."},
  {"role": "user",
   "content": "Generate an SVG icon: outlined credit card with sparkles"},
  {"role": "assistant",
   "content": "<svg xmlns=... viewBox=...><path .../></svg>"}
]}
\end{lstlisting}

\subsection{Data Augmentation}

For each icon we generate up to three training variants with different
user-side prompts pointing to the same SVG target:
(1) the short VLM caption (\textit{"outlined credit card with sparkles"}),
(2) the long VLM caption (full description with style and context), and
(3) the raw icon name from the Iconify metadata (\textit{"credit-card"}).
This teaches the model to handle prompts of varying specificity and phrasing,
reflecting the diversity of real-world user queries.

\subsection{Curriculum Learning}
\label{sec:curriculum}

Following \citet{svgen2025}, training examples are ordered from simple to
complex using a two-level sort key: (1) monochrome icons before multicolor,
(2) ascending path count within each group. This curriculum strategy exposes
the model to the dominant icon patterns first, grounding its SVG priors before
introducing the long-tail of complex multicolor icons.

\subsection{LoRA Fine-tuning}
\label{sec:lora}

We apply Low-Rank Adaptation \citep{hu2022lora} to fine-tune Qwen3.5-9B,
keeping the pre-trained weights frozen and adding trainable low-rank matrices
$A \in \mathbb{R}^{d \times r}$ and $B \in \mathbb{R}^{r \times d}$ at each
attention and MLP projection layer. The effective weight update is:

\[
W' = W + \alpha \cdot A B
\]

where $r$ is the rank (we use $r = 64$) and $\alpha$ is a scaling factor.
With 4-bit quantization of the frozen base weights (QLoRA), total GPU memory
is under 16~GB, enabling training on Apple M4 hardware via the MLX framework
\citep{mlx2023}.

\paragraph{Training configuration.}
\todo{Fill in after training:
batch size, learning rate, number of iterations, hardware, training time per epoch.}
Table~\ref{tab:training_config} summarises the training hyperparameters.

\begin{table}[h]
\centering
\caption{Training configuration for \method{}.}
\label{tab:training_config}
\begin{tabular}{ll}
\toprule
\textbf{Hyperparameter} & \textbf{Value} \\
\midrule
Base model      & Qwen3.5-9B-Instruct \\
Quantization    & 4-bit (QLoRA) \\
LoRA rank $r$   & 64 \\
LoRA $\alpha$   & 128 \\
LoRA target layers & attention + MLP projections \\
Training epochs & 2--3 \\
Batch size      & \todo{X} \\
Learning rate   & \todo{X} \\
Training records & \todo{X} (with augmentation) \\
Hardware        & Apple M4 / \todo{cloud GPU} \\
Training time   & \todo{X hours} \\
\bottomrule
\end{tabular}
\end{table}
